Από τις εξισώσεις του κεφαλαίου \ref{sec:tyre_model} είναι εμφανές πως για τον υπολογισμό θερμοκρασίας και φθοράς στα ελαστικά,
είναι απαραίτητο να γνωρίζουμε τα μεγέθη ολίσθησης καθώς και τα φορτία στους τροχούς.
Ένα υπολογιστικό μοντέλο δεν μπορεί να είναι πιο ακριβές απ'ότι τα δεδομένα εισόδου του, επομένως η καλύτερη προσέγγιση θα ήταν η απ'ευθείας
μέτρηση των μεγεθών ολίσθησης και φορτίων στα ελαστικά. Ωστόσο, οι αισθητήρες μέτρησης ταχυτήτων στους τροχούς (π.χ αισθητήρες doppler) συνδυάζονται από πολύ υψηλό κόστος, ενώ
η μέτρηση των φορτίων στους τροχούς είναι σχεδόν αδύνατη ή καθόλου πρακτική. Τέλος, σε συνθήκες αγώνα, οι περισσότεροι αισθητήρες αφαιρούνται από το όχημα, αφενός για τη μείωση του βάρους αλλά αφετέρου
και για την ελαχιστοποίηση των εξαρτημάτων που μπορούν να αστοχήσουν κατά τη διάρκεια του αγώνα.\\
\noindent
Είναι φανερό, λοιπόν, πως χρειάζεται μια μέθοδος προσέγγισης αυτών των μεγεθών χωρίς την απ'ευθείας μέτρηση τους. Για τον λόγο αυτό αναπτύσσεται μια μεθοδολογία
εκτίμησης της καθολικής κατάστασης του οχήματος με τα ελάχιστα δεδομένα εισόδου.\\[5pt]

Για την εφαρμογή της παρούσας εργασίας, θεωρείται πως είναι γνωστή η ταχύτητα του οχήματος σε κάθε σημείο κατά μήκος της πίστας ενώ επίσης η τροχιά του οχήματος (γεωμετρία της πίστας). Επιπλέον αισθητήρες, όπως για παράδειγμα επιταχυνσιόμετρο στο κέντρο μάζας, μπορούν να αντικαταστήσουν κάποιο βήμα της διαδικασίας αυξάνοντας την ακρίβεια του εκάστοτε τμήματος.
Τέλος, περαιτέρω μεγέθη που χρησιμοποιούνται στη μοντελοποίηση για την εκτίμηση της κατάστασης θεωρούνται γνωστά και σταθερά χαρακτηριστικά του οχήματος που μπορούν να προσδιορισθούν εκ των προτέρων σε εργαστηριακές δοκιμές. 

Η διαδικασία εκτίμησης, παρουσιάζεται σχηματικά στο \cref{fig:slip_est_flow},
και βασίζεται στην αντίστροφη επίλυση του προβλήματος δυναμικής οχήματος. Η επίλυση επιτυγχάνεται αριθμητικά και αναλύεται σε επόμενη παράγραφο.\\
\noindent
Από τις εισόδους του μοντέλου και τις παραδοχές, μπορούν να εκτιμηθούν ή να μετρηθούν οι ταχύτητες του κέντρου μάζας του οχήματος, και οι αντίστοιχες επιταχύνσεις.\\

Όπως θα αναλυθεί στις επόμενες παραγράφους, καταλήγουμε σε δύο ανεξάρτητες μεταβλητές, τη γωνία πλευρικής ολίσθησης $\beta$, και τη γωνία διεύθυνσης των εμπρόσθιων τροχών $\delta$.

% ----- TikZ: high-level slip estimation pipeline -----
\begin{figure}[htbp]
\centering
\begin{tikzpicture}[
    node distance=0.55cm and 1.0cm,
    block/.style   ={rectangle, draw=black, fill=white,      text width=7cm,
                     align=center, rounded corners=4pt, minimum height=0.9cm, font=\small},
    io/.style      ={rectangle, draw=black, fill=black!5,    text width=5.2cm,
                     align=center, rounded corners=2pt, minimum height=0.9cm, font=\small},
    outbox/.style  ={rectangle, draw=black, fill=black!12,   text width=5.2cm,
                     align=center, rounded corners=2pt, minimum height=0.9cm, font=\small\bfseries},
    decision/.style={diamond,   draw=black, fill=white,      text width=2.4cm,
                     align=center, font=\scriptsize, inner sep=1pt},
    arr/.style     ={->, thick, black}
]

  \node[io]                       (inp)    {Είσοδοι: γεωμετρία πίστας, χρονοσειρά ταχύτητας};
  \node[block, below=of inp]   (kinem)  {Υπολογισμός επιταχύνσεων οχήματος $\alpha_X, \alpha_Y$};
  \node[block, below=of kinem]    (startOfIteration)     {Αρχικοποίηση πλευρικής ολίσθησης\\ και γωνίας τιμονιού ($\beta_0, \delta_0$)};
  \node[block, below=of startOfIteration]    (fz)     {Υπολογισμός κατακόρυφων\\ φορτίων ανά τροχό $F_{z,i}$};
  \node[block, below=of fz]    (wheelSlips)     {Υπολογισμός ταχυτήτων και\\ μεγεθών ολίσθησης τροχών $\alpha_i,\kappa_i$};
  \node[block, below=of wheelSlips]    (tyreForces)     {Υπολογισμός πρόσφυσης\\ ελαστικών $F_{u,i}, F_{v,i}$};
  \node[block, below=of tyreForces]       (solve)  {Επίλυση εξισώσεων ισορροπίας\\ (μόνιμη κατάσταση)};
  \node[decision, below=of solve] (conv)   {σύγκλιση;};
  \node[outbox, below=of conv]    (out)    {Έξοδος: ολισθήσεις και δυνάμεις ανά τροχό ($\alpha_i, \kappa_i, F_{u,i}, F_{v,i}$)};

  \draw[arr] (inp)    -- (kinem);
  \draw[arr] (kinem)  -- (startOfIteration);
  \draw[arr] (startOfIteration)     -- (fz);
  \draw[arr] (fz)  -- (wheelSlips);
  \draw[arr] (wheelSlips)     -- (tyreForces);
  \draw[arr] (tyreForces)     -- (solve);
  \draw[arr] (solve)     -- (conv);
  \draw[arr] (conv)   -- node[right, font=\scriptsize]{ναι} (out);

  % failure branch
  \node[font=\scriptsize, right=0.8cm of conv, align=center]
    (fail) {επισήμανση\\ασύγκλιτου σημείου};
  \draw[arr, dashed] (conv.east) -- (fail.west)
      node[midway, above, font=\scriptsize]{όχι};
      
        % ----- dashed box around the iteration loop -----
  \node[draw=black, dashed, rounded corners=6pt,
        inner sep=8pt, fit=(wheelSlips)(fz)(tyreForces)(solve)] 
        (loopbox) {};
        
  \node[anchor=south, font=\normalsize, xshift=-4pt, rotate=90] 
      at (loopbox.west) {Επαναληπτική διαδικασία};

\end{tikzpicture}
\caption{Διάγραμμα ροής -- εκτίμηση ολισθήσεων και δυνάμεων ανά τροχό από δεδομένα γύρου.}
\label{fig:slip_est_flow}
\end{figure}
% ----- τέλος flowchart -----
